The recovered October 1990 working paper by \citet{gentryhoque1990wp}
provides a distinct simultaneous-equation comparison. Its twelve signed
cash-flow components come from income statements and balance-sheet changes.
Each relative component divides by total cash flow. The nine-component
free cash flow after working capital includes investment, dividends and
fixed coverage expenditures as well as operating and working-capital flows
(printed pp.~4--6). Those definitions differ from simply selecting a modern
statement's operating cash-flow line.

Let \(S\) be the sum of those nine normalized components, \(N\) the
normalized financing flow, \(O\) the normalized other-asset/liability flow,
and \(C\) the normalized change in cash and marketable securities. The
paper's equation~(1) then gives
\[
S+N+O-C=0,\qquad S=C-N-O.
\]
The second equality moves financing and other flows to the right; it neither
chooses a financing instrument nor supplies a payment date. Different
financing and cash changes can reconcile the same \(S\) while leaving
different future obligations. Equation~(7) of the retained paper prints its
cash term without the normalization star used in equation~(1). Every term
in the displayed local derivation is explicitly normalized. The source
notation ambiguity remains recorded, and a zero total-cash-flow denominator
has no supplied evaluation rule.

The paper separately proposes five behavioral equations linking operating
flows, receivables, inventory, payables and investment (printed pp.~10--11).
It compares OLS and three-stage least squares using 117 December-year-end
companies. The exhibits report selected significant coefficients; the
discussion describes specification changes based on system fit (printed
pp.~13 and~19). Reported weighted R-squared therefore describes the fitted
sample comparisons. The retained paper supplies neither a held-out forecast
test nor the complete instrument matrix, company data and estimation code
needed for local reproduction. Its interpretation of earlier models remains
the authors' characterization where those models have not been inspected.

For AccDSL, this comparison separates a calculation, a reconciliation and a
decision rule. A definition computes an amount from its inputs; an identity
checks consistency among amounts; a behavioral rule proposes an economic
decision under stated information and assumptions. A simultaneous behavioral
system would also need a solution method, selection among solutions and
explicit failure semantics. The current expression evaluator solves no such
system. The paper instead motivates keeping these requirements distinct from
the dated events and accounting operations that produce financial statements.

